% \section{Safe process with payload}
% \label{sec:safe-process-with}

% The following process is safe and in $\lrProc{n}{\tEmpty \vdash \tyUnit}$.

% \begin{align*}
%   & (\set{},\mset{ 
%     \tLet[\tLVB][\tLVA]{\tApp{\tRecv}{(\tApp{\tFork}{\tAbs{\tLVA}{\tApp{\tClos}{(\tApp{\tSend}{\tPair{\tUnit}{\tLVA}})}}})}}{\tLet*{\tLVB}{\tApp{\tClos}{\tLVA}}}
%      })
%   \\
%   \redu
%   & (\set{ \ch },\mset{ 
%     \tLet[\tLVB][\tLVA]{(\tApp{\tRecv}{\ch})}{\tLet*{\tLVB}{\tApp{\tClos}{\tLVA}}}
%     , 
%     \tApp{(\tAbs{\tLVA}{\tApp{\tClos}{(\tApp{\tSend}{\tPair{\tUnit}{\tLVA}})}})}{\ch}
%      })
%   \\
%   \redu
%   & (\set{ \ch },\mset{ 
%     \tLet[\tLVB][\tLVA]{(\tApp{\tRecv}{\ch})}{\tLet*{\tLVB}{\tApp{\tClos}{\tLVA}}}
%     , 
%     {\tApp{\tClos}{(\tApp{\tSend}{\tPair{\tUnit}{\ch}})}}
%      })
%   \\
%   \redu
%   & (\set{ \ch },\mset{ 
%     \tLet[\tLVB][\tLVA]{\tPair{\tUnit}{\ch}}{\tLet*{\tLVB}{\tApp{\tClos}{\tLVA}}}
%     , 
%     {\tApp{\tClos}{\ch}}
%      })
%   \\
%   \redu
%   & (\set{ \ch },\mset{ 
%     {\tLet*{\tUnit}{\tApp{\tClos}{\ch}}}
%     , 
%     {\tApp{\tClos}{\ch}}
%      })
%   \\
%   \redu
%   & (\set{ \ch },\mset{ 
%     {{\tApp{\tClos}{\ch}}}
%     , 
%     {\tApp{\tClos}{\ch}}
%      })
%   \\
%   \redu
%   & (\set{ },\mset{ 
%     \tUnit
%     , 
%     \tUnit
%      })
% \end{align*}

% We have $
%   \tLet[\tLVB][\tLVA]{\tApp{\tRecv}{(\tApp{\tFork}{\tAbs{\tLVA}{\tApp{\tClos}{(\tApp{\tSend}{\tPair{\tUnit}{\tLVA}})}}})}}{\tLet*{\tLVB}{\tApp{\tClos}{\tLVA}}}
%  \in \lrTerm{n}{\tEmpty \vdash \tyUnit}$ because
% $\tLet[\tLVB][\tLVA]{(\tApp{\tRecv}{c})}{\tLet*{\tLVB}{\tApp{\tClos}{\tLVA}}}
% \in \lrTerm{n}{\ch: \sRecv{\tyUnit}{\tClos} \vdash
%   \tyUnit}$.
% The reduction sequence is similar to the one shown for the first
% example in \cref{sec:safe-process}.

\section{Reductions for leader election example}
\label{sec:reduct-lead-elect}

As an example, we trace the reduction sequence starting from $p_0$  where $\tm_0$
communicates first with $\tm_1$.
\begin{itemize}
\item $\tau$ transition (internal communication) to
  \begin{small}
    \begin{align*}
      \tm_0 &=
                    \tApp{\tClos}{(\tApp{\tSend}{\tPair{\asTm{true}}{\chB}})}
      \\
      \tm_1 &=
                    \tLet[\tLVB][\chB]{\tPair{\asTm{false}}{\chB}}{\tLet*{\tApp{\tClos}{\asTm{cs}}}{\tIf{\tLVB}{\tApp{\tClos}{(\asTm{select}~\asTm{left}~\ch)}}{\tUnit}}}
      \\
      \tm_2 &=
                    \tLet[\tLVB][\chB]{\tApp{\tRecv}{\chB}}{\tLet*{\tApp{\tClos}{\asTm{cs}}}{\tLet*{\tApp{\tClos}{\chB}}{\tIf{\tLVB}{\tApp{\tClos}{(\asTm{select}~\asTm{right}~\ch)}}{\tUnit}}}}
    \end{align*}
  \end{small}
\item $\tau$ transition ($\beta$ pair elimination) to
  \begin{small}
    \begin{align*}
      \tm_0 &=
                    \tApp{\tClos}{(\tApp{\tSend}{\tPair{\asTm{true}}{\chB}})}
      \\
      \tm_1 &=
                    {\tLet*{\tApp{\tClos}{\asTm{cs}}}{\tIf{\asTm{false}}{\tApp{\tClos}{(\asTm{select}~\asTm{left}~\ch)}}{\tUnit}}}
      \\
      \tm_2 &=
                    \tLet[\tLVB][\chB]{\tApp{\tRecv}{\chB}}{\tLet*{\tApp{\tClos}{\asTm{cs}}}{\tLet*{\tApp{\tClos}{\chB}}{\tIf{\tLVB}{\tApp{\tClos}{(\asTm{select}~\asTm{right}~\ch)}}{\tUnit}}}}
    \end{align*}
  \end{small}
\item $\tau$ transition (internal communication) to
  \begin{small}
    \begin{align*}
      \tm_0 &=
                    \tApp{\tClos}{\chB}
      \\
      \tm_1 &=
                    {\tLet*{\tApp{\tClos}{\asTm{cs}}}{\tIf{\asTm{false}}{\tApp{\tClos}{(\asTm{select}~\asTm{left}~\ch)}}{\tUnit}}}
      \\
      \tm_2 &=
                    \tLet[\tLVB][\chB]{\tPair{\asTm{true}}{\chB}}{\tLet*{\tApp{\tClos}{\asTm{cs}}}{\tLet*{\tApp{\tClos}{\chB}}{\tIf{\tLVB}{\tApp{\tClos}{(\asTm{select}~\asTm{right}~\ch)}}{\tUnit}}}}
    \end{align*}
  \end{small}
\item $\tau$ transition ($\beta$ pair elimination) to
  \begin{small}
    \begin{align*}
      \tm_0 &=
                    \tApp{\tClos}{\chB}
      \\
      \tm_1 &=
                    {\tLet*{\tApp{\tClos}{\asTm{cs}}}{\tIf{\asTm{false}}{\tApp{\tClos}{(\asTm{select}~\asTm{left}~\ch)}}{\tUnit}}}
      \\
      \tm_2 &=
                    {\tLet*{\tApp{\tClos}{\asTm{cs}}}{\tLet*{\tApp{\tClos}{\chB}}{\tIf{\asTm{true}}{\tApp{\tClos}{(\asTm{select}~\asTm{right}~\ch)}}{\tUnit}}}}
    \end{align*}
  \end{small}
\item $\tau$ transition (internal communication, closing $\asTm{cs}$,
  new $C = \set{\ch, \chB}$)
  \begin{small}
    \begin{align*}
      \tm_0 &=
                    \tApp{\tClos}{\chB}
      \\
      \tm_1 &=
                    {\tLet*{\tUnit}{\tIf{\asTm{false}}{\tApp{\tClos}{(\asTm{select}~\asTm{left}~\ch)}}{\tUnit}}}
      \\
      \tm_2 &=
                    {\tLet*{\tUnit}{\tLet*{\tApp{\tClos}{\chB}}{\tIf{\asTm{true}}{\tApp{\tClos}{(\asTm{select}~\asTm{right}~\ch)}}{\tUnit}}}}
    \end{align*}
  \end{small}
\item $\tau$ transitions (2x $\beta$ unit)
  \begin{small}
    \begin{align*}
      \tm_0 &=
                    \tApp{\tClos}{\chB}
      \\
      \tm_1 &=
                    {{\tIf{\asTm{false}}{\tApp{\tClos}{(\asTm{select}~\asTm{left}~\ch)}}{\tUnit}}}
      \\
      \tm_2 &=
                    {{\tLet*{\tApp{\tClos}{\chB}}{\tIf{\asTm{true}}{\tApp{\tClos}{(\asTm{select}~\asTm{right}~\ch)}}{\tUnit}}}}
    \end{align*}
  \end{small}
\item $\tau$ transition (internal communication, closing $\chB$,
  new $C = \set{\ch}$)
  \begin{small}
    \begin{align*}
      \tm_0 &=
                    \tUnit
      \\
      \tm_1 &=
                    {{\tIf{\asTm{false}}{\tApp{\tClos}{(\asTm{select}~\asTm{left}~\ch)}}{\tUnit}}}
      \\
      \tm_2 &=
                    {{\tLet*{\tUnit}{\tIf{\asTm{true}}{\tApp{\tClos}{(\asTm{select}~\asTm{right}~\ch)}}{\tUnit}}}}
    \end{align*}
  \end{small}
\item $\tau$ transition (unit elim)
  \begin{small}
    \begin{align*}
      \tm_0 &=
                    \tUnit
      \\
      \tm_1 &=
                    {{\tIf{\asTm{false}}{\tApp{\tClos}{(\asTm{select}~\asTm{left}~\ch)}}{\tUnit}}}
      \\
      \tm_2 &=
                    {{{\tIf{\asTm{true}}{\tApp{\tClos}{(\asTm{select}~\asTm{right}~\ch)}}{\tUnit}}}}
    \end{align*}
  \end{small}
\item $\tau$ transitions (2x $\beta$ bool)
  \begin{small}
    \begin{align*}
      \tm_0 &=
                    \tUnit
      \\
      \tm_1 &=
                    {{{\tUnit}}}
      \\
      \tm_2 &=
                    {{{{\tApp{\tClos}{(\asTm{select}~\asTm{right}~\ch)}}}}}
    \end{align*}
  \end{small}
\item There are two remaining labeled transitions: $\lSend{\ch}{\asTm{right}}$
  followed by $\lTerm{\ch}$. The resulting process is terminated and is
  an element of $\lrProc{n}{\tEmpty \vdash \tyUnit}$:
  \begin{small}
    \begin{align*}
      \tm_0 &=
                    \tUnit
      \\
      \tm_1 &=
                    {{{\tUnit}}}
      \\
      \tm_2 &=
                    {{{{\tUnit}}}}
    \end{align*}
  \end{small}
\end{itemize}
\section{Duality and Equivalence of Recursive Session Types}
\label{sec:app-duality}

\begin{figure}[tp]
  \textbf{Session types} $\sty$, extending \cref{sec:extens-recurs-ii}
  \begin{align*}
    \sty &::= \sRecv{\ty}{\sty} \sepr \sSend{\ty}{\sty} \sepr
                   \sTerm \sepr \tyMu\tvar.\sty \sepr
                   \tvar \sepr \dual{\tvar}
    % & \text{($\dual{\tvar}$ never occurs in a source type)}
  \end{align*}
  \textbf{Duality} $\dual{\sty}$
  \begin{align*}
    \dual{\sRecv{\ty}{\sty}} &= \sSend{\ty}{\dual{\sty}}
    &
    \dual{\sSend{\ty}{\sty}} &= \sRecv{\ty}{\dual{\sty}}
    &
    \dual{\sTerm} &= \sTerm
    &
    \dual{\tyMu\tvar.\sty} &= \tyMu\tvar.\,\sflipop_{\tvar}(\dual{\sty})
    \span\span\span
  \end{align*}
  Duality exchanges \emph{every} variable $\tvar$ with 
  $\dual{\tvar}$; $\sflipop_\alpha$ does the same swap, but only for the
  designated variable $\tvar$, leaving every other variable alone. 

  \textbf{Equivalence} $\ty \approx \tyB, \sty \approx \styB$, coinductively
  \begin{mathpar}
    \begin{bussproof}[Eq-Unit]
      \bussAx{\tyUnit \approx \tyUnit}      
    \end{bussproof}
    \and
    \begin{bussproof}
      \bussAs{\ty \approx \ty'}
      \bussAs{\tyB \approx \tyB'}
      \bussB{\lFunc{\ty}{\tyB} \approx {\lFunc{\ty'}{\tyB'}}}
    \end{bussproof}
    \and
    \begin{bussproof}
      \bussAs{\ty \approx \ty'}
      \bussAs{\tyB \approx \tyB'}
      \bussB{\lPair{\ty}{\tyB} \approx {\lPair{\ty'}{\tyB'}}}
    \end{bussproof}
    \\
    \begin{bussproof}[Eq-Recv]
      \bussAs{\ty \approx \tyB}
      \bussAs{\sty \approx \styB}
      \bussB{\sRecv{\ty}{\sty} \approx \sSend{\tyB}{\styB}}
    \end{bussproof}
    \and
    \begin{bussproof}[Eq-Send]
      \bussAs{\ty \approx \tyB}
      \bussAs{\sty \approx \styB}
      \bussB{\sSend{\ty}{\sty} \approx \sRecv{\tyB}{\styB}}
    \end{bussproof}
    \and
    \begin{bussproof}[Eq-Close]
      \bussAx{\sTerm \approx \sTerm}
    \end{bussproof}
    \\
    \begin{bussproof}[Eq-MuL]
      \bussAs{\sty[\tyMu\tvar.\sty/\tvar] \approx \styB}
      \bussU{\tyMu\tvar.\sty \approx \styB}
    \end{bussproof}
    \and
    \begin{bussproof}[Eq-MuR]
      \bussAs{\sty \approx \styB[\tyMu\tvar.\styB/\tvar]}
      \bussU{\sty \approx \tyMu\tvar.\styB}
    \end{bussproof}
  \end{mathpar}

  \textbf{Relational duality} $\sty \isdual \styB$, coinductively
  \begin{mathpar}
    \begin{bussproof}[D-Recv]
      \bussAs{\ty \approx \tyB}
      \bussAs{\sty \isdual \styB}
      \bussB{\sRecv{\ty}{\sty} \isdual \sSend{\tyB}{\styB}}
    \end{bussproof}
    \and
    \begin{bussproof}[D-Send]
      \bussAs{\ty \approx \tyB}
      \bussAs{\sty \isdual \styB}
      \bussB{\sSend{\ty}{\sty} \isdual \sRecv{\tyB}{\styB}}
    \end{bussproof}
    \and
    \begin{bussproof}[D-Close]
      \bussAx{\sTerm \isdual \sTerm}
    \end{bussproof}
    \\
    \begin{bussproof}[D-MuL]
      \bussAs{\sty[\tyMu\tvar.\sty/\tvar] \isdual \styB}
      \bussU{\tyMu\tvar.\sty \isdual \styB}
    \end{bussproof}
    \and
    \begin{bussproof}[D-MuR]
      \bussAs{\sty \isdual \styB[\tyMu\tvar.\styB/\tvar]}
      \bussU{\sty \isdual \tyMu\tvar.\styB}
    \end{bussproof}
  \end{mathpar}
  % \ie $\sty \isdual \styB$ holds if $\unf(\sty)$ and
  % $\unf(\styB)$ match one of the four conclusions, up to the
  % symmetry of \ruleLabel{D-Var}.

  \textbf{Facts}
  \begin{align*}
    \dual{\dual{\sty}} &= \sty
    &
    \dual{\unf(\sty)} &= \unf(\dual{\sty})
    &
    \sty \isdual \styB &\iff \sty \approx \dual{\styB}
  \end{align*}
  \caption{Equivalence and duality of recursive session types
    \quad \rocq{dual}}
  \Description{The grammar of recursive session types with polarized
    variables, the duality function including its case for the
    recursive binder, the eight coinductive rules for type equivalence, the five
    rules of coinductive relational duality, and the three facts relating the
    two presentations.}
  \label{fig:duality}
\end{figure}

\Cref{fig:duality} provides the definitions for
\cref{sec:extens-recurs-ii}. Duality
\cite{LindleyM16} is a function on types, and on a recursive
type it cannot simply descend into the body: under $\tyMu\tvar.\sty$ the
binder stands for the dual protocol, so occurrences of $\tvar$
that were meant to denote the original type one must change polarity. This is why the
grammar carries a dualized variable $\dual{\tvar}$ (\rocq{SVarD}) and
why the $\mu$ case applies $\sflipop_\alpha$ (\rocq{sflip_ty}) to the dualized
body: tail occurrences, which continue as the session itself, come back as
$\tvar$, while occurrences inside a payload keep the polarity they were
given. Taking $\dual{\tvar} = \tvar$ instead is unsound, and
the type of a channel end would no longer be the dual of the type of the other
end.

Duality also has a relational reading $\sty \isdual \styB$. The two agree:
$\sty \isdual \styB$ holds iff
$\sty \approx \dual{\styB}$ (\rocq{sdual_seq_dual})
\cite{DBLP:journals/corr/abs-2004-01322}. 

\section{Reduction Semantics and Harmony}
\label{sec:app-harmony}
Our operational semantics is a labeled transition system, and
\emph{reduction} is defined as its $\tau$-fragment
(\cref{sec:calculus}). This appendix shows that nothing hinges on
that presentation choice: we define a conventional, self-contained
reduction semantics on processes --- call-by-value evaluation
contexts, $\beta$-redexes, and one rule per interaction, with no
labels anywhere --- and prove that it coincides with $\tau$-reduction
(\cref{thm:harmony}).

\paragraph{Evaluation contexts}
Evaluation contexts extend the contexts of
\cref{fig:term-semantics} by  including all extensions of \cref{sec:extensions}.
\begin{align*}
  \ectx ::={} & \hole
  \sepr \tApp{\ectx}{\tm}
  \sepr \tApp{\val}{\ectx}
  \sepr \tPair{\ectx}{\tm}
  \sepr \tPair{\val}{\ectx}
  \sepr \tLet*{\ectx}{\tm} \\
  & \sepr \tLet[\vr][\vrB]{\ectx}{\tm}
  \sepr \tInj{1}{\ectx}
  \sepr \tInj{2}{\ectx}
  \sepr \tCase{\ectx}{\ivx_1}{\tm_1}{\ivx_2}{\tm_2}
\end{align*}

\paragraph{Pure reduction}
The label-free term reduction
$\tm \redu[\mathsf{V}] \tm'$ consists of the
$\beta$-redexes, closed under evaluation contexts. It collects the term
reductions of the base calculus and the extensions.
\begin{align*}
  \tApp{(\tAbs{\ivx}{\tm})}{\val} & \redu[\mathsf{V}] \tSub{\tm}{\val}{\ivx}
  & \tLet*{\tUnit}{\tm} & \redu[\mathsf{V}] \tm
  \\
  \tLet[\vr][\vrB]{\tPair{\valB}{\val}}{\tm} & \redu[\mathsf{V}] \tSub{\tSub{\tm}{\valB}{\vr}}{\val}{\vrB}
  & \tCase{(\tInj{\asTm{i}}{\val})}{\ivx_1}{\tm_1}{\ivx_2}{\tm_2} & \redu[\mathsf{V}] \tSub{\tm_{\asTm{i}}}{\val}{\ivx_{\asTm{i}}}
  \\
  \tApp{\tForce}{\IExp{(\ivx = \tm)}} & \redu[\mathsf{V}] \tSub{\tm}{\IExp{(\ivx = \tm)}}{\ivx}
  & \ectx[\tm] & \redu[\mathsf{V}] \ectx[\tm'] \quad \text{if } \tm \redu[\mathsf{V}] \tm'
\end{align*}

\begin{lemma}[\rocq{vred_reduction}]
  \label{lem:vred-reduction}
  $\tm \redu[\mathsf{V}] \tm'$
  if and only if
  $\tm \trans{\beta} \tm'$.
\end{lemma}

\paragraph{Reduction semantics for processes}
\Cref{fig:red-semantics} gives the reduction semantics
$\proc \sred \proc'$: a thread takes a pure step
(\ruleLabel{R-Lift}); a fork spawns a thread and restricts a fresh
channel (\ruleLabel{R-Fork}); a send and a receive on the same
restricted channel synchronize (\ruleLabel{R-Com}); two closes on the
same restricted channel retire it (\ruleLabel{R-Close}); and a
forwarder fuses a restricted end with its partner by substitution
(\ruleLabel{R-Link}, and symmetrically for the other end of the
pair). Interactions happen under evaluation contexts in the
participating threads; where
the LTS of \cref{fig:proc-semantics} lets the term transitions
surface the interacting redexes, these rules find them
syntactically.

\pt{TODO: one direction is simple; what about the converse?}
This reduction semantics is still different from the semantics of
\citet{conf/esop/LindleyM15}, which works on structured process terms. However,
it is easy to relate structured process terms to a formulation based
on a soup of processes like ours.

\begin{figure}[tp]
  \textbf{Reduction semantics} $\proc \sred \proc'$
  \begin{mathpar}
    \begin{bussproof}[R-Lift]
      \bussAs{\tm \redu[\mathsf{V}] \tm'}
      \bussU{(\chset, \pool \uplus \mset{\tm}) \sred (\chset, \pool \uplus \mset{\tm'})}
    \end{bussproof}
    \and
    \begin{bussproof}[R-Fork]
      \bussAs{\ch \text{ fresh}}
      \bussU{(\chset, \pool \uplus \mset{\ectx[\tApp{\tFork}{\val}]}) \sred
             (\chset \cupdot \set{\ch}, \pool \uplus \mset{\ectx[\ch], \tApp{\val}{\ch}})}
    \end{bussproof}
    \and
    \begin{bussproof}[R-Com]
      \bussAs{\ch \in \chset}
      \bussU{(\chset, \pool \uplus \mset{\ectx[\tApp{\tSend}{\tPair{\val}{\ch}}], \ectx'[\tApp{\tRecv}{\ch}]}) \sred
             (\chset, \pool \uplus \mset{\ectx[\ch], \ectx'[\tPair{\val}{\ch}]})}
    \end{bussproof}
    \and
    \begin{bussproof}[R-Close]
      \bussAs{\ch \in \chset}
      \bussU{(\chset, \pool \uplus \mset{\ectx[\tApp{\tClos}{\ch}], \ectx'[\tApp{\tClos}{\ch}]}) \sred
             (\chset \setminus \set{\ch}, \pool \uplus \mset{\ectx[\tUnit], \ectx'[\tUnit]})}
    \end{bussproof}
    \and
    \begin{bussproof}[R-Link]
      \bussAs{\ch \in \chset}
      \bussU{(\chset, \pool \uplus \mset{\ectx[\tApp{\tLink}{\tPair{\ch}{\chB}}]}) \sred
             (\chset \setminus \set{\ch}, \tSub{(\pool \uplus \mset{\ectx[\tUnit]})}{\chB}{\ch})}
    \end{bussproof}
  \end{mathpar}
  \caption{Reduction semantics: pure steps, fork, communication,
    closing, and forwarding, with no labels \quad \rocq{red}}
  \Description{Five reduction rules for process soups: lifting a pure
    term step, forking with a fresh channel, communication between a
    send and a receive redex, closing a channel from both ends, and
    forwarding by substitution.}
  \label{fig:red-semantics}
\end{figure}

\paragraph{Harmony}
Both directions of \cref{thm:harmony} rest on the correspondence
between the congruence closure of the labeled semantics and
evaluation contexts: filling a labeled step into a context is a
labeled step, and conversely every labeled step decomposes into an
evaluation context around a head step.
\ruleLabel{R-Lift} then matches
\ruleLabel{P-Act-Beta} by \cref{lem:vred-reduction}, and each
interaction rule matches the process rule that consumes the
corresponding labels: \ruleLabel{R-Fork} with
\ruleLabel{P-Act-Fork}, \ruleLabel{R-Com} with \ruleLabel{P-Act-Msg},
\ruleLabel{R-Close} with \ruleLabel{P-Act-Close}, and
\ruleLabel{R-Link} with \ruleLabel{P-Act-Link}, whose unordered label
pair accounts for the two orientations of the rule.

%%% Local Variables:
%%% mode: latex
%%% TeX-master: "main"
%%% End:
